\section{The cubic trace and its exact denominator}\label{sec:arithmetic}
For the numerical calculation let $F_g$ and $E_g$ be the disconnected
degree-generating series at scalar and opposite weights. Set $D_n=D_{n+1,3}$
and $n=g-1$. The transition to the integral lattice will use the full
row of coefficients of $D_n$, including coefficients in genera greater
than the genus of the base.
\subsection{The quadratic and cubic corrections}

Write $a_{g,m}=X_{g,m}(-1)$ and $x_{g,m}=X_{g,m}(1)$.
With $s=2\sin(u/2)$ the scalar and antidiagonal degree-two residues
are, respectively,
\[
 s^{4n}\big((4-2s)^n+(4+2s)^n\big),\qquad
 2s^{4n}(4-s^2)^n,\qquad n=g-1
\]
by \cite[\S8]{BP}. The linear terms of the two series agree. Thus
the difference $J_n=[q^2](\log F_{n+1}-\log E_{n+1})$ is
\begin{equation}\label{eq:arith-two}
 J_n(w)=2\sum_{l=1}^n4^{n-l}
 \left(\binom n{2l}-(-1)^l\binom nl\right)w^{2n+l}.
\end{equation}
In particular $J_1=2w^3$, $J_2=24w^5-2w^6$. The first formula and
Theorem~\ref{thm:main} give the universal identity
\[
 a_{2,2}=x_{2,2}-2x_{4,1}.
\]
At this genus the only nonconstant modes of $X_{2,2}(z)$ have
weights $\pm1$ and are equal. Each is $x_{4,1}/2$, as in the
geometric double-cover calculation of \cite{WueProduct}.

At degree three let $w^3\lambda_i$ be the three eigenvalues of the
balanced scalar handle operator. Its characteristic equation after
division by $w^3$ in the eigenvalue variable is
\begin{equation}\label{eq:arith-char}
 X^3-AX^2+\tfrac13 BX-\tfrac13 C=0,
\end{equation}
where, with $L_0=w^2+54w-243$,
\[
 A=w^2+50w+81,\quad B=4(w-6)L_0,\quad
 C=4(w-4)(w-3)^2L_0.
\]
This is the scalar specialization of the partition multiplication and
cap operators of \cite[\S6 and Appendix]{BP}, derived in Appendix~\ref{app:handle}. The antidiagonal eigenvalues are
$U,U,V$, where $U=w^3(4-w)(3-w)^2$ and $V=w^3(3-w)^2$ by the
hook-length formula. Writing $T_n=\sum_i\lambda_i^n$ gives
\begin{equation}\label{eq:arith-difference}
 D_n=w^{3n}T_n-2U^n-V^n-w^n J_n.
\end{equation}
Here $[q]F_{n+1}=[q]E_{n+1}=w^n$. The final term removes the
contribution of a degree-one and a degree-two connected component.
For example $T_1=A$ and $T_2=A^2-2B/3$. Substitution gives
\begin{align}
 D_1&=120w^4-20w^5+2w^6,\notag\\
 D_2&=14448w^7-1136w^8+\tfrac{4720}{3}w^9
             -332w^{10}+40w^{11}-2w^{12}.\label{eq:arith-row}
\end{align}
Let $\iota(w^{h-1})=x_{h,1}$, extended linearly. Recovering the
disconnected universal generators from their logarithms gives
\[
 a_{3,3}=x_{3,3}-\iota(D_2)
          +x_{3,1}(-24x_{6,1}+2x_{7,1}).
\]
Indeed the quadratic difference here is $-\iota(J_2)$, and the
linear generators agree. In particular the coefficient of $x_{10,1}$
in the scalar-minus-antidiagonal cubic logarithm is $4720/3$. The corresponding
coefficient in the disconnected scalar residue before the logarithm
and the change of weights is $[w^9]w^6T_2=292/3$.

The two positive odd weights contributing to $[w^9]D_2$ can also
be distinguished. Write $b_1,b_3$ for their respective coefficients
at $x_{10,1}$. All higher odd weights vanish at this genus by
Theorem~\ref{thm:main}. Let $B_d(z)$ denote the balanced handle at
weights $(z,1)$ in degree $d$, with the normalizations in
Appendix~\ref{app:handle}. The genus-three connected coefficient is
\[
 \operatorname{tr}(B_3(z)^2)-w^2\operatorname{tr}(B_2(z)^2)+w^6/3.
\]
For clarity, its values at the two additional weight ratios needed to
separate the modes are
\[
\begin{array}{c|cc}
 z&2&-2\\ \hline
 [w^9]&549&-19891/6 .
\end{array}
\]
They follow by substituting $(t_1,t_2)=(\pm2,1)$ in the multiplication,
metric and cap in Appendix~\ref{app:handle} and reducing with
$r^2-(2-w)r+1=0$, $r=e^{iu}$. Equality of the opposite weights and
subtraction at $z=1,-1$ and $z=2,-2$ give
\[
 4(b_1+b_3)=4720/3,\qquad 5b_1+\tfrac{65}{4}b_3=23185/6.
\]
Thus $b_1=674/3$ and $b_3=506/3$ in the universal algebra.

\subsection{The valuation of the cubic trace}

Write $v_3$ for the usual valuation on rational numbers, with
$v_3(3)=1$. The common denominator of a finite list of rational
numbers is its least positive integral multiplier making the entire
list integral.

\begin{theorem}\label{thm:arith-rows}
The common denominator of the coefficients of $D_n$ is $3^{s_n}$,
where
\[
 s_{2k}=k,\qquad s_{2k+1}=\max\{0,k-2-v_3(2k+1)\}.
\]
For $s_n>0$ the polynomial $3^{s_n}D_n$ has the following reduction
in $\mathbb F_3[w]$:
\begin{align*}
 n=2k:&\qquad 2(-1)^kw^{9k},\\
 n=2k+1:&\qquad
 2(-1)^k\frac{2k+1}{3^{v_3(2k+1)}}w^{9k+2}(w-1).
\end{align*}
\end{theorem}
\begin{proof}
Equation~\eqref{eq:arith-char} gives a recurrence over
$\ZZ[1/3][w]$ for $T_n$, with $T_0=3$, $T_1=A$ and
$T_2=A^2-2B/3$. In \eqref{eq:arith-difference} all terms other than
$w^{3n}T_n$ are integral polynomials. Work in the completion of
$\QQ(w)$ for the Gauss valuation at three, with residue field
$\mathbb F_3(w)$. Its valuation on a polynomial is the least valuation
of a coefficient. The polynomial
\[
 P(X)=3X^3-3AX^2+BX-C
\]
reduces to a linear polynomial with derivative $\bar B=w^3$.
Hensel's lemma gives a unit root $z_0$ with
$\bar z_0=\bar C/\bar B=\bar A$. For the two remaining roots let
$S=A-z_0$ be their sum and $T=C/(3z_0)$ their product. Since
\[
 P(A)=AB-C=72L_0(3w^2-14w-25),
\]
the divided difference $(P(A)-P(z_0))/(A-z_0)$ reduces to $\bar B$,
and hence
\[
 v_3(S)=2,\quad \overline{S/9}=2(w-1)/w,\qquad
 v_3(T)=-1,\quad \overline{3T}=w^3.
\]
The remaining roots have equal valuation $-1/2$ in a splitting
field, since their sum has positive valuation and their product
has valuation $-1$. Their ratio reduces to $-1$. For $k>0$ their
$2k$th power sum divided by $(-T)^k$ therefore reduces to $2$.
Adding the integral term $z_0^{2k}$ and multiplying by
$3^kw^{6k}$ proves the even case. Directly $D_0=0$.

For $n=2k+1$ the power sum of the remaining pair is
\[
 \sum_{l=0}^k\frac{n}{2l+1}\binom{k+l}{2l}(-T)^{k-l}S^{2l+1}.
\]
This is the coefficient of $t^n$ in
$-t\frac{d}{dt}\log(1-St+Tt^2)$. The $l=0$ term has valuation
$v_3(n)-k+2$. Each subsequent term has valuation greater by at least
$5l-v_3(2l+1)>0$. If the first valuation is nonnegative, the entire
trace is integral at three. Otherwise clearing its denominator leaves
only the first term in the reduction. Inserting the reductions of
$S/9$ and $3T$ and the factor $w^{3(2k+1)}$ gives the second formula.
\end{proof}

\begin{proposition}
Only primes at most $m$ divide $\operatorname{den}(x_{g,m})$.
For a prime $p$ and $g\ge2$ the $p$-part of
$\operatorname{den}(x_{g,p})$ is the $p$-part of the common
denominator of the coefficients of $D_{g,p}(w)$.
\end{proposition}

\begin{proof}
In the balanced split bundle use scalar weights $(1,1)$. The
elementary symmetric functions of the $d$ projected fiber coordinates on each
of the two normal lines are sections of their tensor powers of
orders at most $d$. Their evaluations at finitely many points
detect every coefficient, and simultaneous vanishing means that
the cycle is the $d$-fold zero section. The supported integral
Thom class is divided by a product of weights from $1$ through $d$
in Pardon's construction \cite[Definition~4.7]{Pardon}.
Hence $x_{g,d}$ is integral at every prime $p>d$.

In the difference of the degree-$p$ logarithms, the only term
with denominator $p$ from the logarithm uses $p$ copies of the common
linear coefficient. It cancels. Every term involving between two
and $p-1$ factors belongs to
$H_{(p)}=H_{\ZZ}\otimes\ZZ_{(p)}$: the denominator from the logarithm is a
unit at $p$ and the indices of its factors are smaller than $p$.
Also $a_{g,p}$ is integral. Theorem~\ref{thm:main} gives
$x_{g,p}-\iota(D_{g,p})\in H_{(p)}$. Because the $x_{h,1}$ belong
to the polynomial basis, clearing the denominator at $p$ on their
span is equivalent to clearing it on their coefficients.
\end{proof}

\begin{proof}[Proof of the denominator formula in Theorem~\ref{thm:cubic}]
For genera zero and one the local models of Section~\ref{sec:statement}
have empty separating tuple: the two lines of $\cO(-1)^{\oplus2}$
have no nonzero sections after any positive-degree cover, and in genus
one one can take a nontorsion degree-zero line and its inverse. A
section of a pulled-back degree-zero line would trivialize it and
taking the norm would make a positive power of the chosen line
trivial. Thus the whole compact cycle space in each multiple is the
point, independent of weights. For $g\ge2$ take $n=g-1$.
Theorem~\ref{thm:main} at degree two and equation~\eqref{eq:arith-two}
give $x_{g,2}-a_{g,2}=\iota(J_n)$ with $J_n\in\ZZ[w]$.
Thus every scalar double is integral. Expanding
the cubic logarithms gives
\[
 x_{g,3}=a_{g,3}+x_{g,1}\iota(J_n)+\iota(D_{g,3}).
\]
The first two terms on the right are integral, so the denominator
over $\ZZ$ is exactly the common coefficient denominator of
$D_{g,3}$. Its value is $3^{s_n}$ by the valuation in Theorem~\ref{thm:arith-rows}. In particular no factor of two
survives.
\end{proof}


Theorem~\ref{thm:arith-rows} also identifies the entire leading residue: in $3^{s_n}x_{n+1,3}\bmod3H_{\ZZ}$,
replace each $w^{h-1}$ in that theorem by $x_{h,1}$ whenever $s_n>0$.
For $g=2,3,4,5,6,7,8$ the resulting denominators are
$1,3,1,9,1,27,3$.

For arbitrary nonzero normal weights $(r,s)$ in the genus-two
quadratic case, the symmetry and the Laurent bound give
$X_{2,2}(r/s)=Y+(r/s+s/r)Z$. The difference at equal and opposite
weights is $2x_{4,1}$, whence
$X_{2,2}(r/s)=x_{2,2}+(r-s)^2x_{4,1}/(2rs)$.

For comparison, at a
genus-three core, the scalar-minus-opposite-weight connected logarithm
has coefficient $4720/3$ at $(2\sin(u/2))^{18}$. At fixed arithmetic
source genus ten its coefficient at $u^{18}$ is $23593/5$. Finally,
for a general stable normal bundle a neighborhood with a general second
jet isolates the triple and its point class exceeds the scalar class
by $(8/3)x_{10,1}$. Here eight counts cyclic covers and three is the
weight of their amplitude.
